\documentclass[9pt,a4paper]{article}

\usepackage[a4paper,landscape,margin=0.4cm]{geometry}
\usepackage{multicol}
\usepackage{amsmath,amssymb,bm}
\usepackage{mdframed}
\usepackage{tikz}
\usepackage{titlesec}
\usepackage{url}
\usepackage{hyperref}
\usepackage{graphicx}
\setlength{\columnsep}{0.35cm}
\setlength{\parindent}{0pt}
\setlength{\parskip}{2pt}

\titleformat{\section}{\large\scshape}{}{0pt}{}
\titleformat{\subsection}{\normalsize\bfseries}{}{0pt}{}
\titlespacing*{\section}{0pt}{2pt}{2pt}
\titlespacing*{\subsection}{0pt}{1pt}{1pt}

\begin{document}
% % ======= ST2132-style 3-part dividers =======
% \begin{tikzpicture}[remember picture,overlay]
% % Line at 1/3 width
% \draw[black!60, line width=0.8pt]
%   ($(current page.south west)!0.333!(current page.south east)$) --
%   ($(current page.north west)!0.333!(current page.north east)$);

% % Line at 2/3 width
% \draw[black!60, line width=0.8pt]
%   ($(current page.south west)!0.666!(current page.south east)$) --
%   ($(current page.north west)!0.666!(current page.north east)$);
% \end{tikzpicture}

\scriptsize

\begin{center}
\textbf{\Large CN3104 Cheatsheet AY25/26}\\
Github:\href{https://github.com/Shaunteojiwu?tab=repositories}{Shaunteojiwu}\\
\href{https://www.linkedin.com/in/shaun-teo-jiwu0824/}{LinkedIn} \\
\end{center}

\begin{multicols}{3}

% ==========================================
% NONLINEAR EQUATIONS
% ==========================================
\section{Nonlinear Equations}

Goal: Solve $f(x)=0$.

\subsection{Graphical Method}
-Plot $f(x)$ vs $x$; roots where curve crosses x-axis.\\
-Simple but imprecise.

\subsection{Incremental Search}
-Using computer's brute force (relies on trying all possibilities)\\
1) Divide $[x_{\min},x_{\max}]$ into $n$ intervals.\\
2) Find sign changes: if $f(x_i)f(x_{i+1})<0$, a root lies inside.\\
-Use MATLAB function 'incsearch'\\
-Default number of intervals, ns=50

\textbf{Disadvantages:}  
-Misses roots if step too large; slow if step too small.

\subsection{Bracketing Method:Bisection}
1) Guess 2 initial starting point, where the root is within.\\
$f(a)f(b)<0$.(1 positive,1 negative-can apply incremental search/graphical method for initial guess).\\  
2) Iterate midpoint:
\[c=\frac{a+b}{2}.\]
3) If $f(a)f(c)<0$ set $b=c$, else $a=c$.
4) Stop when stopping criterion met. (threshold error, similar to error in fixed-point iteration)\\
\textbf{Guaranteed convergence}, slow.

\subsection{Fixed-Point Iteration}
1) Rewrite $x=g(x)$.\\
2) Iterate $x_{i+1}=g(x_i)$.\\
3) Calculate error,
Error:
\[
e_a=\left|\frac{x_{i+1}-x_i}{x_{i+1}}\right|100\%.
\]
-Stop when error within 0.01
-For bisection, error xi+1 is new root for the present iteration

%Converges if $|g'(x)|<1$.

\subsection{Newton-Raphson}
% \begin{figure}
%     \centering
%     \includegraphics[width=0.5\linewidth]{newtonraphson.png}
%     \caption{Enter Caption}
%     \label{fig:placeholder}
% \end{figure}
\begin{center}
\includegraphics[width=0.7\linewidth]{newtonraphsonfull.jpg}
\end{center}
% \[
% x_{i+1}=x_i-\frac{f(x_i)}{f'(x_i)}.
% \]
-Fast convergence. \\
-Requires derivative and good initial guess.\\
-Conditions for stop same as secant method.

\subsection{Secant Method}
No derivative needed, uses backward finite difference to estimate f'(x):
\[
x_{i+1}=x_i - f(x_i)\frac{x_i-x_{i-1}}{f(x_i)-f(x_{i-1})}.
\]
-Requires 2 starting value (do not need to bracket root)\\
-Conditions for stop:
$x_i \approx constant$ or $f(xi) \approx 0$


% ==========================================
% LINEAR SYSTEMS
% ==========================================
\section{Systems of Linear Equations}

Solve $A\mathbf{x}=\mathbf{b}$.

\subsection*{Gaussian Elimination vs.\ LU Factorization}

\textbf{Gaussian Elimination (GE):}
Goal: convert $A \to U$ (upper triangular).\
Uses row operations: 
    \[
    R_2 \leftarrow R_2 - \frac{a_{21}}{a_{11}} R_1.
    \]\
Matrix $A$ is modified during elimination.\
Does not store multipliers.\

\textbf{LU Factorization:}
    Goal: write $A = LU$.\
    $U$ is obtained from elimination (same as GE).\
    $L$ stores the multipliers used in elimination.\
    To avoid changing $A$, start $L$ as the identity matrix and insert multipliers.\


\textbf{Identity Matrix Trick:}
\[
L = I 
\quad\Rightarrow\quad
L(2,1) = m_{21} = \frac{a_{21}}{a_{11}}
\]
Insert each multiplier into $L$ while forming $U$.

\textbf{Example:}
\[
A = \begin{bmatrix} 2 & 4 \\ 6 & 5 \end{bmatrix},
\qquad
m_{21} = \frac{6}{2} = 3
\]
\[
U = \begin{bmatrix} 2 & 4 \\ 0 & -7 \end{bmatrix},
\qquad
L = \begin{bmatrix} 1 & 0 \\ 3 & 1 \end{bmatrix}
\]
\[
A = LU
\]

Forward elimination to make $U$ upper triangular.  
Then back-substitute.

Pivot equation: eliminate variables below pivot using  
\[
\text{factor}=\frac{a_{ik}}{a_{kk}}.
\]

\subsection{Ill-Conditioning}
Small changes in data cause large changesthe  in solution.

% ==========================================
% LINEAR REGRESSION
% ==========================================
\section{Linear Regression}

Given data $(x_i,y_i)$.

\subsection{Simple Linear Regression}
{Goal:} Fit a straight line that best predicts $y$ from $x$\
A model is linear if it is linear in the coefficients (parameters).\
Model:  
\[
y = a_0 + a_1 x + e.
\]

Normal equations:
\[
a_0 n + a_1\sum x_i=\sum y_i
\]
\[
a_0\sum x_i + a_1\sum x_i^2=\sum x_iy_i.
\]

\subsection{General Least Squares}
Matrix form:  
\[
\hat{y}=Xb
\]
\[
\varepsilon = y - Xb
\]
Minimize $\Phi=\varepsilon^T\varepsilon$.

Solution:
\[
b = (X^T X)^{-1} X^T y.
\]

% ==========================================
% INTEGRATION
% ==========================================
\section{Numerical Integration}

Goal:
\[
I=\int_a^b f(x)\,dx.
\]

\subsection{Trapezoidal Rule}
Single interval:
\[
I\approx\frac{b-a}{2}[f(a)+f(b)].
\]

Composite:
\[
I=\frac{h}{2}[f(x_0)+2\sum_{i=1}^{n-1}f(x_i)+f(x_n)].
\]

Error:
\[
E_a=-\frac{(b-a)^3}{12n^2}f''(\xi).
\]

\subsection{Simpson's 1/3 Rule}
Requires even number of intervals:\
\[
I = (b - a)\,\frac{f(x_0) + 4f(x_1) + f(x_2)}{6}
\]
Composite:\
\[
I = (b - a)\,
\frac{
    f(x_0)
    + 4 \sum_{\substack{i = 1 \\ i \text{ odd}}}^{n-1} f(x_i)
    + 2 \sum_{\substack{j = 2 \\ j \text{ even}}}^{\,n-2} f(x_j)
    + f(x_n)
}{3n}
\]

\subsection{Simpson's 3/8 Rule}
\[
I = (b - a)\,\frac{f(x_0) + 3f(x_1) + 3f(x_2) + f(x_3)}{8}
\]


% ==========================================
% DIFFERENTIATION
% ==========================================
\section{Numerical Differentiation}

Using finite differences.

\subsection{Forward Difference}
\[
f'(x_i)\approx\frac{f(x_{i+1})-f(x_i)}{h}.
\]
Error $O(h)$.

\subsection{Backward Difference}
\[
f'(x_i)\approx\frac{f(x_i)-f(x_{i-1})}{h}.
\]
Error $O(h)$.

\subsection{Central Difference}
\[
f'(x_i)\approx\frac{f(x_{i+1})-f(x_{i-1})}{2h}.
\]
Error $O(h^2)$.

% ==========================================
% Ordinary Differential Equations
% ==========================================
\section{ODEs}
Solve:
\[
\frac{dy}{dt} = f(t,y),\quad y(t_0)=y_0.
\]

\subsection{Euler Method}
Use finite difference approximation\
\[
\frac{dv}{dt} \approx \frac{\Delta v}{\Delta t}
= \frac{v(t_{i+1}) - v(t_i)}{t_{i+1} - t_i}
\]
Rearranging the expression\
\[
v_{i+1}=v_i + \frac{dvi}{dt}{\Delta t}
\]\
General form:\
\[
y_{i+1} = y_i + f(t_i, y_i)\, h
\]



\subsection{Heun's Method (Improved Euler)}
\[
y_1 = y(t_0) 
    + \frac{h}{2} \Big( 
        f(t_0, y_0) 
        + f\big(t_1,\, y_0 + h f(t_0, y_0)\big)
      \Big)
\]

Predictor:
\[
k_1=f(t_i,y_i)
\]
Corrector:
\[
k_2=f(t_i+h,y_i+h k_1)
\]
Update:
\[
y_{i+1}=y_i+\frac{h}{2}(k_1+k_2).
\]

\subsection{Midpoint Method}
\[
y_{i+1} = y_i + f\!\left(t_{i+\frac{1}{2}},\, y_{i+\frac{1}{2}}\right) h
\]

\[
k_1=f(t_i,y_i)
\]
\[
k_m=f(t_i+\tfrac{h}{2},y_i+\tfrac{h}{2}k_1)
\]
\[
y_{i+1}=y_i+h k_m.
\]

\subsection{Runge-Kutta 4 (RK4)}
\[
y_{i+1}
= y_i
+ \frac{1}{6}\bigl(k_1 + 2k_2 + 2k_3 + k_4\bigr)h
\]
\[
k_1 = f(t_i,\, y_i)
\]

\[
k_2 = f\!\left(t_i + \frac{h}{2},\, y_i + \frac{h}{2}k_1\right)
\]

\[
k_3 = f\!\left(t_i + \frac{h}{2},\, y_i + \frac{h}{2}k_2\right)
\]

\[
k_4 = f(t_i + h,\, y_i + hk_3)
\]

% \[
% k_1=f(t,y)
% \]
% \[
% k_2=f(t+h/2, y + h k_1/2)
% \]
% \[
% k_3=f(t+h/2, y + h k_2/2)
% \]
% \[
% k_4=f(t+h, y + h k_3)
% \]
% \[
% y_{i+1}=y_i+\frac{h}{6}(k_1+2k_2+2k_3+k_4).
% \]

\subsection{Adaptive RK (ode45)}
MATLAB uses RK4(5) pair.  
Small step when function changes quickly; large step otherwise.

% ==========================================




% \section*{Matrix \& Vector Rules (MATLAB Style)}

% \subsection*{Basic Shapes}
% \[
% y \in \mathbb{R}^{n\times 1},\quad 
% b \in \mathbb{R}^{m\times 1},\quad
% X \in \mathbb{R}^{n\times m}
% \]

% \subsection*{Transpose Rules}
% \[
% (A + B)^T = A^T + B^T
% \]
% \[
% (A - B)^T = A^T - B^T
% \]
% \[
% (AB)^T = B^T A^T
% \]

% \subsection*{Inner Products \& Norms}
% \[
% a^T b = \text{scalar}
% \]
% \[
% \|a\|^2 = a^T a = \sum_i a_i^2
% \]
% \[
% \text{SSE} = (y - Xb)^T (y - Xb)
% \]

% \subsection*{Useful MATLAB Equivalents}
% \begin{verbatim}
% a' * b      % inner product
% a' * a      % squared norm
% e  = y - X*b;
% SSE = e' * e;
% \end{verbatim}

% \subsection*{Derivative Rules}
% \[
% \frac{\partial}{\partial b}(c^T b) = c
% \]
% \[
% \frac{\partial}{\partial b}\left(b^T A b\right)
% = (A + A^T)b
% \]
% If \(A\) is symmetric:
% \[
% \frac{\partial}{\partial b}(b^T A b) = 2Ab
% \]

% \subsection*{Least Squares Gradient}
% \[
% \Phi(b) = (y - Xb)^T (y - Xb)
% \]
% \[
% \nabla_b \Phi = 2X^T(Xb - y)
% \]

% \subsection*{Normal Equations}
% \[
% X^T X\, b = X^T y
% \]
% \[
% b = (X^T X)^{-1}X^T y
% \]

% \subsection*{MATLAB Implementation}
% \begin{verbatim}
% e    = y - X*b;
% grad = -2 * X' * e;

% b_hat = (X' * X) \ (X' * y);
% \end{verbatim}
% \begin{multicols}{2}
% \small

% % ====== TITLE ======
% \begin{center}
% \textbf{\Large Matrix \& Vector Rules (MATLAB Style)}
% \end{center}

% % ====== BOX 1: SHAPES ======
% \noindent\fbox{
% \parbox{\linewidth}{
% \textbf{Basic Shapes}
% \[
% y \in \mathbb{R}^{n\times 1},\quad 
% b \in \mathbb{R}^{m\times 1},\quad
% X \in \mathbb{R}^{n\times m}
% \]
% }}

% \vspace{5pt}

% % ====== BOX 2: TRANSPOSE ======
% \noindent\fbox{
% \parbox{\linewidth}{
% \textbf{Transpose Rules}
% \[
% (A+B)^T = A^T + B^T
% \]
% \[
% (A-B)^T = A^T - B^T
% \]
% \[
% (AB)^T = B^T A^T
% \]
% }}

% \vspace{5pt}

% % ====== BOX 3: NORMS & INNER PRODUCTS ======
% \noindent\fbox{
% \parbox{\linewidth}{
% \textbf{Inner Products \& Norms}
% \[
% a^T b = \text{scalar}, \quad
% \|a\|^2 = a^T a
% \]
% \[
% \text{SSE} = (y - Xb)^T (y - Xb)
% \]
% \texttt{
% a' * b \quad \% inner product \\
% a' * a \quad \% squared norm
% }
% }}

% \vspace{5pt}

% % ====== BOX 4: DERIVATIVES ======
% \noindent\fbox{
% \parbox{\linewidth}{
% \textbf{Derivative Rules}
% \[
% \frac{\partial}{\partial b}(c^T b) = c
% \]
% \[
% \frac{\partial}{\partial b}(b^T A b) 
% = (A + A^T)b
% \]
% If \(A\) is symmetric:
% \[
% \frac{\partial}{\partial b}(b^T A b) = 2Ab
% \]
% }}

% \vspace{5pt}

% % ====== BOX 5: LEAST SQUARES ======
% \noindent\fbox{
% \parbox{\linewidth}{
% \textbf{Least Squares}
% \[
% \Phi(b) = (y - Xb)^T (y - Xb)
% \]
% \[
% \nabla_b \Phi = 2X^T(Xb - y)
% \]
% \[
% X^T X\, b = X^T y
% \]
% \[
% b = (X^T X)^{-1}X^T y
% \]
% \texttt{
% e = y - X*b; \\
% grad = -2 * X' * e; \\
% b\_hat = (X' * X) \ (X' * y);
% }
% }}

% \end{multicols}
% \begin{multicols}{2}
% \small

% \begin{center}
% \textbf{\Large Matrix \& Vector Rules (MATLAB Style)}
% \end{center}

% % ========== BOX 1 ==========
% \noindent\fbox{
% \parbox{\linewidth}{
% \textbf{Basic Shapes}
% \[
% y\!:\!n\!\times\!1,\quad
% b\!:\!m\!\times\!1,\quad
% X\!:\!n\!\times\!m
% \]
% }}

% \vspace{4pt}

% % ========== BOX 2 ==========
% \noindent\fbox{
% \parbox{\linewidth}{
% \textbf{Transpose Rules}
% \[
% (A+B)^T=A^T+B^T
% \]
% \[
% (A-B)^T=A^T-B^T
% \]
% \[
% (AB)^T = B^T A^T
% \]
% }}

% \vspace{4pt}

% % ========== BOX 3 ==========
% \noindent\fbox{
% \parbox{\linewidth}{
% \textbf{Inner Products \& Norms}
% \[
% a^T b = \text{scalar},\qquad
% \|a\|^2 = a^T a
% \]
% \[
% \text{SSE} = (y - Xb)^T(y - Xb)
% \]

% {\ttfamily\small
% a' * b   % inner product\\
% a' * a   % norm squared
% }
% }}

% \vspace{4pt}

% % ========== BOX 4 ==========
% \noindent\fbox{
% \parbox{\linewidth}{
% \textbf{Derivative Rules}
% \[
% \frac{\partial}{\partial b}(c^T b) = c
% \]
% \[
% \frac{\partial}{\partial b}(b^T A b) 
% = (A + A^T)b
% \]
% If \(A\) symmetric:
% \[
% \frac{\partial}{\partial b}(b^T A b)=2Ab
% \]
% }}

% \vspace{4pt}

% % ========== BOX 5 ==========
% \noindent\fbox{
% \parbox{\linewidth}{
% \textbf{Least Squares}
% \[
% \Phi(b)=(y-Xb)^T(y-Xb)
% \]
% \[
% \nabla_b\Phi
% =2X^T(Xb-y)
% \]
% \[
% X^TX\,b = X^T y
% \]
% \[
% b = (X^TX)^{-1}X^T y
% \]

% {\ttfamily\small
% e = y - X*b;\\
% grad = -2 * X' * e;\\
% b\_hat = (X' * X) \ (X' * y);
% }
% }}
% \begin{multicols}{2}
% \small

% \begin{center}
% \textbf{\Large Matrix \& Vector Rules}
% \end{center}

% \begin{mdframed}[linewidth=0.5pt]
% \textbf{Basic Shapes}\\[-4pt]
% \[
% y\in\mathbb{R}^{n\times1},\ 
% b\in\mathbb{R}^{m\times1},\ 
% X\in\mathbb{R}^{n\times m}
% \]
% \end{mdframed}

% \begin{mdframed}[linewidth=0.5pt]
% \textbf{Transpose Rules}\\[-4pt]
% \[
% (A+B)^T=A^T+B^T
% \]
% \[
% (A-B)^T=A^T-B^T
% \]
% \[
% (AB)^T=B^T A^T
% \]
% \end{mdframed}

% \begin{mdframed}[linewidth=0.5pt]
% \textbf{Inner Products \& Norms}\\[-4pt]
% \[
% a^T b=\text{scalar},\quad 
% \|a\|^2=a^T a
% \]
% \[
% \text{SSE}=(y-Xb)^T(y-Xb)
% \]

% \ttfamily\small
% a' * b   % inner product\\
% a' * a   % squared norm
% \end{mdframed}

% \begin{mdframed}[linewidth=0.5pt]
% \textbf{Derivative Rules}\\[-4pt]
% \[
% \frac{\partial}{\partial b}(c^Tb)=c
% \]
% \[
% \frac{\partial}{\partial b}(b^TAb)=(A+A^T)b
% \]
% If \(A\) symmetric:
% \[
% \frac{\partial}{\partial b}(b^TAb)=2Ab
% \]
% \end{mdframed}

% \begin{mdframed}[linewidth=0.5pt]
% \textbf{Least Squares}\\[-4pt]
% \[
% \Phi(b)=(y-Xb)^T(y-Xb)
% \]
% \[
% \nabla_b\Phi
% =2X^T(Xb-y)
% \]
% \[
% X^TX\,b=X^T y
% \]
% \[
% b=(X^TX)^{-1}X^T y
% \]

% \ttfamily\small
% e = y - X*b;\\
% grad = -2 * X' * e;\\
% b\_hat = (X' * X) \ (X' * y);
% \end{mdframed}


\end{multicols}
\end{document}
